% ============================================================
%  Modelo de nota de aula — copie com ./nova-nota.sh
%
%  Compile com:  make pdf/<nome>.pdf
%  Tema claro :  make claro
% ============================================================
\documentclass{notaaula}        % use [claro] para fundo branco

\titulonota{Título da aula}
\creditos{3º ano · Física · Mês/Ano · Prof. Josué Cavalcante}

\begin{document}
\cabecalho

% Resumo de uma ou duas linhas: conteúdo, turmas, contexto.
\begin{sobre}
Assunto · subtópicos · turmas · escola.
\end{sobre}

% ============================================================
\section{Primeiro tópico}

% Caixa "COPIAR": o que o aluno leva para o caderno.
\begin{copiar}{Nome curto do bloco}
\definicao Texto da definição. Use \dest{negrito} nas palavras-chave.
\[ y = a x + b \]
\simbolos{$y$ (unidade); $x$ (unidade).}
\diaadia{Uma frase ligando o conceito ao cotidiano.}
\end{copiar}

% Outros rótulos coloridos disponíveis dentro das caixas:
%   \definicao  \formulas  \relacoes  \modelo  \resolucao
%   \rotulo{Texto livre.}   -> rótulo azul qualquer
%   \atencao{...}           -> alerta âmbar
%   \resultado{...}         -> destaque coral (respostas)

% Figura com legenda (TikZ, pgfplots ou \includegraphics).
\begin{figuranota}{Legenda explicando o que a figura mostra.}
\begin{tikzpicture}
  \draw[naAzul, thick, -{Stealth}] (0,0) -- (3,0) node[right,
        font=\tiny, color=naSuave] {$x$};
\end{tikzpicture}
\end{figuranota}

% ============================================================
\section{Segundo tópico}

Texto corrido fora de caixa, para o que não precisa ser copiado.
\begin{itens}
\item primeiro item;
\item segundo item.
\end{itens}

% Exemplo resolvido (caixa verde).
\begin{exemplo}
Enunciado do exemplo.
\begin{alternativas}
\item primeira pergunta;
\item segunda pergunta.
\end{alternativas}
\resolucao (a) desenvolvimento. \resultado{Resposta em destaque.}

(b) desenvolvimento.
\end{exemplo}

% Dica / erro comum (caixa âmbar).
\begin{dica}
O engano que a turma costuma cometer, e como evitar.
\end{dica}

% ============================================================
% Lista de exercícios. A numeração das questões é contínua
% entre os três níveis.
\begin{exercicios}
\nivel{1}{Conceitos}
\begin{questoes}
\item Questão de múltipla escolha.
  \begin{alternativas}
  \item alternativa a;
  \item alternativa b;
  \item alternativa c;
  \item alternativa d.
  \end{alternativas}
\item Questão aberta.
\end{questoes}

\nivel{2}{Aplicação}
\begin{questoes}
\item Questão de aplicação.
\end{questoes}

\nivel{3}{Mais trabalhados}
\begin{questoes}
\item Questão mais longa.
\end{questoes}

\gabarito{1b · 2) resposta · 3) resposta · 4) resposta.}
\end{exercicios}

\end{document}
